\section{Complexity-theoretic preliminaries}
\label{sec:prelim}

This section introduces the machinery Section~\ref{sec:findings-tables} uses. A
reader fluent in game-based security definitions and restricted algorithm
classes can move on.

\subsection{Games, advantage, and the baseline}

A security definition in cryptography is a game between a challenger and an
adversary. The challenger samples a secret, the adversary interacts under stated
rules, and the adversary wins when it produces something the definition names.
The adversary's \emph{advantage} is its winning probability minus the
probability it wins without the object under attack.

Membership inference instantiates this. The challenger trains a model on a
dataset that either contains a target record or does not, the adversary sees the
model, and it guesses which. The challenger flips a fair coin, so an adversary
that ignores the model wins half the time, and one half is the baseline.

The baseline belongs to the definition and not to the evaluation that follows
it. Yeom et al.\ show what a careless choice produces. Under their membership
advantage the two worlds are equiprobable before the adversary consults the
model, so the marginal quantity is zero by construction and only model access
produces advantage \cite{YGFJ18}. Balle, Cherubin and Hayes give the
constructive form for reconstruction. Their baseline is the best attack that
ignores the model entirely, and their criterion carries a tolerance $\eta$
instead of demanding exact recovery \cite{BCH22}.
Section~\ref{sec:findings-tables} records, for every work, whether it states a
baseline of either kind.

\subsection{Adversary classes}

A security claim is a statement about some set of algorithms. Six such sets
appear in this corpus, and Section~\ref{sec:findings-tables} keeps them apart.

\emph{Polynomial time} is the default in cryptography. The adversary is a
probabilistic algorithm whose running time is polynomial in a security
parameter, and a scheme is secure when every such algorithm has negligible
advantage. A claim is \emph{unconditional} when it holds outright and
\emph{conditional} when it holds under an assumption such as the hardness of
factoring, of learning parity with noise, or of planted clique.

\emph{Circuit size} is the non-uniform version. The adversary is a family of
Boolean circuits, one per input length, and the resource is the number of gates.
Non-uniform adversaries are at least as strong as uniform ones, so a bound
against circuits is the stronger claim.

\emph{Statistical queries} restrict how an algorithm touches data. An SQ
algorithm never sees a sample. It submits a bounded function and receives that
function's expectation over the distribution, accurate to a stated tolerance.
Gradient descent and principal component analysis both fit the model, so SQ
lower bounds carry weight in machine learning \cite{FGRVX17}.

\emph{Low-degree polynomials} restrict an algorithm to thresholding a polynomial
of degree at most $D$ in the observations \cite{Wei25}. Under conditions the
low-degree and statistical-query classes agree \cite{BBHLS21}.

\emph{Bounded memory} restricts space instead of time. Raz shows that learning
parity from a stream requires either quadratic memory in the dimension or
exponentially many samples \cite{Raz16}, and the bounded-storage line in
cryptography builds on the same restriction.

\emph{Function classes} restrict the adversary to a hypothesis class instead of
a resource. Chaudhuri, Imola and Machanavajjhala bound an adversary by the class
its attack algorithm is drawn from, using restricted $f$-divergences
\cite{CIM19}.

Two of these classes are often treated as interchangeable with running time.
Buhai, Hsieh, Jain and Kothari disprove the quasi-polynomial low-degree
conjecture in the regime it was stated. They exhibit planted distributions with
vanishing low-degree advantage whose noisy distinguishing problem is solved far
faster than the conjecture allows \cite{BHJK25}. Their polynomial-time
distinguisher covers a rectangular generalisation, and in the square symmetric
regime the running time is $n^{O(\log^{1/(k-1)} n)}$. Hsieh et al.\ establish
the positive counterpart \cite{HKKLMT26}. A claim resting on a low-degree bound
therefore asserts something different from a claim resting on a time bound.

Hardness moves between problems by reduction, and three things stay behind. A
reduction from problem $A$ to problem $B$ turns an algorithm for $B$ into one
for $A$, so hardness of $A$ transfers to $B$ \cite{BB20}. It does not carry
across adversary classes, so a reduction sound against polynomial-time
algorithms asserts nothing about statistical-query ones. It preserves tightness
only when the author charges the transformation's own cost. And it inherits
every antecedent of the result it starts from.

\subsection{The gap between two parties}

Fine-grained cryptography asks for security against adversaries that exceed the
honest party's resources by a stated factor instead of superpolynomially. Three
conventions for that factor appear in this corpus and they are not
interchangeable.

Brzuska and Couteau work in non-uniform circuit size, and their
$(\varepsilon,\delta)$-fine-grained one-way function admits no inverter of size
below $|f|^{1+\delta}$ \cite{BC22}. LaVigne, Lincoln and Vassilevska Williams
work in uniform word-RAM time and name the object, defining the relative gap as
the constant $\delta$ separating an honest cost of $T(n)^{1-\delta}$ from a
threshold of $T(n)$ \cite{LLW19}. Ball, Rosen, Sabin and Vasudevan also work in
word-RAM time, and they index their error by the adversary's own discount as
well as by the input length, so one number does not describe their notion
\cite{BRSV17}. Degwekar, Vaikuntanathan and Vasudevan take a fourth route,
restricting both parties to a circuit class instead of relating their costs by
an exponent \cite{DVV16}.

One ceiling on such a gap is often quoted without its scope. Barak and Mahmoody
show that an adversary making $O((n/\delta)^2)$ \emph{oracle queries} breaks any
$n$-query key exchange in the random oracle model, and they note that their
adversary runs efficiently in a relativized world where $\mathrm{P} =
\mathrm{NP}$ \cite{BM09}. The bound forbids a super-quadratic query gap in an
idealized model. It constrains neither time gaps nor gaps built on structured
assumptions. In the privacy-preserving machine learning literature the same
word, fine-grained, names per-user access control, so a keyword search returns
that sense and not this one.
